% Pursuing Stacks (cotes 134-2 to 134-8): archive cote 134-3, site page 1 (PDF page 2), typescript page 1.
% typescript. Draft compared: ps2.tex:1-67.
% AI-transcribed by Claude Opus 5.5 (claude-opus-5-5, adaptive thinking, effort high), Message Batches API, batch ps1 (msgbatch_01Rv5hCaNTtUt4JHx1yqoHPL), 2026-11-12; not yet checked by eye.
\scanpage{134-3/1}
\margin{1}
\mspage{1}

Reflections on homotopical algebra

27.2.83.

\margin{\unsure{\#}} The following notes are the continuation of the reflection star-%
ted in my letter to Daniel Quillen written previous week (19.2 - 23.2),
which I will cite by (L) (``letter''). I begin with some corrections and
comments to this letter.

\hand{14}. Homotopical algebra can be viewed as being concerned mainly with
the study of spaces of continu\xout{\unsure{o}}ous \xout{functions} \inserted{maps between spaces} \deleted{and their}\inserted{, and the} algebraic analo-%
gons, \inserted{of these} \xout{\illegible{1}} with a special emphasis on homotopies between such \xout{functions}\inserted[typed]{maps},
homotopies between homotopies etc. The kind of restrictive properties
imposed on \deleted{functions} \inserted{the maps under consideration} is exemplified by the typical example when deman-%
ding that the \xout{functions} \deleted{considered} \inserted[typed]{maps} should be extensions, or liftings, of
a given \xout{function} map. \deleted{It} \inserted{Homotopical algebra} is not \inserted{directly} suited for the study of spaces of homeo-%
morphisms, spaces of immersions, of embeddings, of fibrations etc
- and it seems that the \deleted{latter} study \inserted{of such spaces} has not really yet taken off the
ground. Maybe the main obstacle here lies in the wildness phenomena,
which however, one feels, makes a wholly artificial difficulty, stem-%
ming from the particular way by which topological intuition has been
mathematically formalized, in terms of the basic notion of topological
spaces and continuous maps between them. This transcription, while
adequate for the homotopical point of view, and partly adequate too
for the use of analysts, is rather coarsely unad\hand{e}quate, it seems to me,
in most other geometrico-topological contexts, and particularly so when
i\over{\unsure{l}}{t} comes to studying spaces of homeomorphisms, immersions etc (in all
those questions when \xout{the} ``isotopy'' is replacing the rather coarse homo-%
topy relation), as well as for a study of stratified structures, when
it becomes undispensable to give intrinsic and precise meaning to such
notions as tubular neighbourhoods etc. For a structure theory of strati
fications, it turns out,\inserted[typed]{(somewhat surprisingly maybe)} that even the somewhat cumbersome context of
\xout{pre}topoi and \unsure{pro-topoi|pretopoi} is better suited than topological spaces, and mo\unsure{re}
over directly applicabl\unsure{e} to unconventional contexts such as étale topo-%
logy of schemes, where the conventional transcription of topological
intuition in terms of topological spaces is quite evidently breaking
down. To emphasize one point I was making in (L,p.1), it seems to me
that this breakdown is almost as evident in \xout{\illegible{1}} isotopy questions\hand{,} or
for the needs of a structure theory of stratified ``spaces'' (whatever we
mean by ``spaces'' ...). \deleted{and} \over{i}{I}t is a matter of amazement to me that this
breakdown has not been clearly noticed, and \inserted{still less} overcome by \xout{axsuitable}
working out a more suitable transcription of topological objects and
